% 第一章 佟家的奴才
% 原书页码：003；标题
\section{佟家的奴才}
\par
% 原书页码：003；段落
顺治十二年二月二十七日（1655年4月3日），耶稣会传教士利%
类思和安文思给顺治皇帝呈上六件礼物外加一份中文写的奏本，感谢%
顺治皇帝以及清朝长期以来对他们的优容和恩养。在奏本中，他们简%
要介绍了他们一路走来的经历。奏本是这样写的：%
\par
% 原书页码：003；段落
{\bookkai 大西洋耶稣会士远臣利类思、臣安文思谨奏：}%
\par
% 原书页码：003；段落
{\bookkai 为微臣感恩无地，敬进西国方物，以表臣忱事。臣等产于西洋，}%
{\bookkai 从幼弃家修道。明季东来至蜀，居堂传教，历有年所，不意遭寇掳待}%
{\bookkai 戮。于顺治三年，幸逢大兵西剿，恭遇肃王恩释，依随来京。仰荷皇}%
{\bookkai 上隆恩，特送礼部及光禄寺，图赖固山家豢养，朝夕焚祝万寿，经今}%
{\bookkai 五载，感激无涯。兹有方物六种，躬献阙庭，伏祈皇上，俯赐敕收。\NoteRef{1}{1}}%
\par
% 原书页码：003–004；段落（跨页续段）
两位神父明朝末年就到了中国。利类思是西西里人（现属于意大利），%
1636年（崇祯九年）来华；安文思是葡萄牙人，1642年到成都，和%
利类思一起在那里传教。后来二人一直一起在中国传教，直到安文思%
1677年去世。%
\par
% 原书页码：004；段落
当1655年二人上这份奏本的时候，他们已在中国待了快二十年，%
早已不是来华的新人了。奏本中的遣词造句也可以看出，二人对于清%
初官场政治很了解。奏本感谢了肃王豪格和佟图赖固山（固山是八旗%
内的官职，是一旗的统帅，后改为都统）两个人，然后还提到了两个%
机构：礼部和光禄寺。为什么会感谢人的同时，又感谢机构，而不是%
感谢负责这个机构的人呢？这恰恰显示出了二人的圆滑老练，因为他%
们要隐藏和一个人的关系。%
\par
% 原书页码：004；段落
这个人就是摄政王多尔衮（1612—1650）。当二人在1648年跟着%
豪格的征西大军班师归来抵京时，他们的身份有两个：一是叛贼张献%
忠朝廷的逆臣；二是战地俘虏，已归为奴籍。\NoteRef{1}{2}所以二人到京之时，是%
戴罪之身。当时在京做官的同为耶稣会神父的德国人汤若望，既不敢%
和他们联系也不敢把他们收留在他的教堂中。\NoteRef{1}{3}但当时实际掌控北京%
的多尔衮，并没有为难二人，反而把二人交给专管接待外交使节的礼%
部来照顾，并且还给予了他们在北京城中活动的自由。\NoteRef{1}{4}对于在成都%
张献忠的大西朝当过官的人来说，这二人最初能在北京没有被追究，%
并平稳地安顿下来，不得不说多尔衮对二人是有恩的。但是，二人上%
奏本的时候，多尔衮已经去世五年，早已不是一个可以在朝堂上提及%
的人。多尔衮的爵位在他死后几个月，就因为朝堂政治斗争被剥夺了。\NoteRef{1}{5}%
两个神父很清楚当时的政治气氛，知道哪些名字可以提，哪些不能提。%
\par
% 原书页码：004–005；段落（跨页续段）
当然，他们提到的两个人，豪格和佟图赖，都是当时政治气氛中%
对于他们有利的人。那两个神父又到底跟这两位满人权贵有什么联%
系？这其实就涉及两位神父是如何从张献忠朝廷的罪臣身份一步步和%
满人权贵搭上线的问题。从二人1648年进京算起，到1677年安文思%
去世、1682年利类思去世，他们在满人统治下的北京待了超过三十年，%
历经了两位清朝皇帝。要了解二人怎么慢慢进人满人朝廷的故事，就%
得从他们与两位满人权贵的关系说起。%
% 原书页码：005；标题
\subsection{战场上的俘虏}
\par
% 原书页码：005；段落
跟二位神父相关的大清历史大概得追溯到1643年，也即满人人关%
的前一年。后世称为清太宗文皇帝的皇太极（1592—1643），在没有%
留下遗诏指定继承人的情况下，这一年突然驾崩了。这不可避免地引%
发了各方势力的争位行动。简单来说，当时有两方势力不分伯仲，一%
派支持皇太极的长子豪格继位，另一派支持皇太极的弟弟多尔衮。为%
了避免满人内部的分裂，双方妥协达成了第三个方案：推举皇太极的%
第九子、当时六岁的福临，也就是后来的顺治皇帝登位。作为附加条%
件，多尔衮和济尔哈朗（努尔哈赤之侄）一起为顺治当摄政王。虽然%
最后皇位问题得以妥协解决，但是争位的多尔衮和豪格叔侄双方成了%
仇敌。%
\par
% 原书页码：005–006；段落（跨页续段）
1648年3月，也就是满人控制北京的第四年，多尔衮以谋反的罪%
名逮捕了豪格。逮捕前一个月，豪格刚结束四川西征，剿灭了张献忠，%
得胜回京。豪格在被逮捕当月就死于狱中。无论正史、野史都基本一%
致认为豪格是被多尔衮害死的。最让时人惊诧的是，多尔衮在豪格死%
后，还把豪格的夫人纳人府中。就连北京的汤若望神父也把这当成奇%
谈写进了他发回欧洲的信中。%
\par
% 原书页码：006；段落
1650年12月，多尔衮去世。在他死后两个月，顺治皇帝就在济%
尔哈朗的支持下剥夺了多尔衮家的爵位，瓦解了多尔衮在朝廷中的残%
余势力。同时，豪格的所有爵位得到了恢复。1651年以后，多尔衮就%
不是适合在朝堂上攀附的名字了。%
\par
% 原书页码：006；段落
有了这个背景，才可以理解为什么两位神父在奏本中提到了豪格，%
而不提多尔衮。不过，两位神父提豪格也不能算是攀附。因为二人确%
实和豪格有实在的联系。豪格是西征张献忠大军的最高统帅，而两位%
神父于1644至1647三年间在张献忠成都的朝廷中做官。\NoteRef{1}{6}按他们自己%
的话说，他们是张献忠大军进人成都后，为了活命不得不听命于张献%
忠。他们从西医角度，认为张献忠患有精神疾病，也亲眼见证了张献%
忠残酷掠杀老百姓。后世各种关于张献忠暴虐的传闻，都能在二人的%
记录中得到印证。%
\par
% 原书页码：006；段落
两位神父是在战场上遭遇豪格大军的，两人都在战场上中箭倒地。%
他们的欧洲人面孔救了他们，赶上来补刀的满人士兵没有杀他们，反%
而把他们带到豪格的大帐中。所幸当时随军的满汉语翻译和北京的汤%
若望神父有私交，救了二人的命。有了随军翻译的这层关系，二人虽%
然是战俘，并没有被为难，还得到了特别的关照。\NoteRef{1}{7}%
\par
% 原书页码：006–007；段落（跨页续段）
从四川被俘到最后1648年落脚北京，二人在豪格部队中跟着满%
人士兵一起走了一年半左右。\NoteRef{1}{8}很多满人士兵也知道二人，知道豪格%
关照他们。有一次，豪格传令大部队，让士兵在各自俘虏中寻找两位%
神父的一个仆人。最初，抓住这个仆人的满人士兵并不想把他交出来，%
后来听说是豪格在找这个人，这个士兵亲自把仆人带到豪格帐内，当%
成礼物送给了豪格。\NoteRef{1}{9}%
\par
% 原书页码：007；段落
满人传统中，士兵行军打仗都是自备武器和粮草，擒获的战俘归%
士兵私有，成为士兵的奴隶。对于士兵来说，获得战利品和获取奴隶%
都是同等重要的经济回报。在满人世界中，奴隶和牛、马、衣服这些%
一样，都是私有财产。奴隶最大的来源就是战场上的俘虏，所以俘获%
仆人的士兵一开始并不想交人；后来他把仆人交给豪格，和献给豪格%
一匹布是差不多的。%
\par
% 原书页码：007；段落
由于有豪格的照顾，满人士兵不久也把两个神父当成了自己人。%
大概在二人被俘五个月以后，部队回京路过西安。这时，两个神父已%
经知道满人军队内部的行事习俗。当时满人军队进人城市，按例都要%
临时征用寺庙等公共设施来驻扎。两位神父率先找到了当地的天主教%
堂，住进去后，就对后来的士兵吹嘘这里已经被部队内某某大人物征%
用了。士兵听他们这样说，就离开另寻他处去了。\NoteRef{1}{10}%
\par
% 原书页码：007–008；段落（跨页续段）
当然，两位神父跟着满人部队，适应满人生活的过程还是很痛苦%
的。从他们第一天在战场遇到满人士兵，他们就知道这些士兵是“外%
国人”。这些人和他们熟悉的汉人不一样，不会说汉语。尽管二位神%
父早已精通汉语，无奈从士兵到豪格，都不通汉语，完全没办法交流，%
他们所有交流都需要翻译。\NoteRef{1}{11}1满人军人吃的也和汉人不一样，两位神%
父得跟着满人学他们在行军中吃马肉。\NoteRef{1}{12}跟着满人军队以后，两位神%
父的基本身份也跟以前在明朝以及在张献忠大西朝廷的不一样了，因%
为现在他们不算皇帝的臣民，而就是擒获他们的人的奴隶，是他们主%
子的私有财产。最开始的时候，两位神父各自有自己的主子。这样他%
们二人的生活，比如吃什么，吃多少，都由二人的主子决定。利类思%
运气较好，他的主子比较大方，能基本吃饱。但安文思就没有那么好%
运气，他的主子连基本食物都不给够，他好几次都差点饿死在路上。\NoteRef{1}{13}%
\par
% 原书页码：008；段落
在行军路上的某个时候（1647—1648），二位神父的奴籍被转到了%
豪格名下。主奴关系是满人八旗世界中最基本的阶层关系之一。在满%
人的理念中，主子给奴才提供生活所需，而奴才对主子尽忠。\NoteRef{1}{14}满人%
统治者也深知奴隶是重要的经济资源。因而他们推崇主奴之间以家为%
单位，和睦相处。用努尔哈赤的话说，就是要形成一种“奴才爱主子，%
主子爱奴才”的社会。\NoteRef{1}{15}由于主奴关系中蕴含着家庭关系，因而外人%
在对待一个奴才的时候得看主人是谁；同样，要是主人升官了，自家%
奴才也跟着高兴，感觉脸上有光。%
\par
% 原书页码：008；段落
满人人关以后，汉族大臣极力反对满人的主奴统治，但满人上层%
还是力排众议，坚持主奴关系是满人八旗制度的基础。最有名的例子%
就是清初严厉执行《逃人法》，这个法令的核心就是重罚那些胆敢容%
留逃跑奴隶的人。\NoteRef{1}{16}也是因为这个原因，当汤若望神父知道两位神父%
是奴隶身份，就不敢在他设在北京的教堂中收留他们了。\NoteRef{1}{17}%
\par
% 原书页码：008–009；段落（跨页续段）
对于做过豪格奴才的这一身份，两位神父并没有隐瞒。因为满人%
社会中的奴才和汉人社会的“奴”是不一样的。汉人的“奴”属于贱%
民阶层，地位在贵族、官吏、平民之下。简单来说，在汉人社会，“奴”%
和其他属于贱民阶层的人社会地位低，在社会多方面受歧视，而且%
法律和社会习俗还对贱民阶层有很多专门的行为约束。\NoteRef{1}{18}但在满人社%
会，奴才的身份不代表一个人的社会地位，他们的社会地位主要取%
决于他们主子的社会地位。用通俗的话来说，在满人世界，重要的不%
是你是否是奴才身份，而是做的是谁家的奴才。如果是皇帝和达官贵%
人家的奴才，那是一种有利的身份。%
\par
% 原书页码：009；段落
两位传教士知道做过豪格奴才，表明了他们算是豪格一脉的人。%
这个背景对二人在顺治朝廷是有利的，因而他们在给顺治上奏本的时%
候专门提到这个已经过世了七年的前主子。%
% 原书页码：009；标题
\subsection{康熙母亲的娘家}
\par
% 原书页码：009；段落
在豪格1648年去世以后，两位神父的奴籍转到了佟图赖名下。\NoteRef{1}{19}%
佟图赖是正蓝旗汉军都统，是整个佟氏家族的核心人物。八旗每一旗%
有一位旗主，旗下又分为满、蒙、汉三旗，也就有三位都统。由于旗%
主必须是努尔哈赤家族的后代，所以都统这个职位是非努尔哈赤后代%
在八旗中能达到的最高职位。%
\par
% 原书页码：009；段落
明末以来，佟氏家族一直是满人龙兴之地——辽东地区的名门望%
族。努尔哈赤娶的第一位夫人就来自佟家，他还用“佟”作为他的汉%
族姓氏。佟氏家族从努尔哈赤时期起就在满人政坛中有重要地位。满%
人1644年人关以后，佟家子弟遍布朝野，当时甚至有“佟半朝”一说。%
\par
% 原书页码：009–010；段落（跨页续段）
在满人政坛中，佟氏家族的最大特点是其成员兼具满人和汉人两%
方面的文化特点。由于康熙的母亲出自佟家，因而佟家是汉人还是满%
人就决定了康熙到底有没有汉人血统的问题。历史上一直争论不清到%
底他们是汉化了的满人还是满人化了的汉人。\NoteRef{1}{20}°但无论观点如何，都%
承认佟氏兼具了汉满两方面的特质。在17世纪初，明清在争夺辽东%
的时候，两方也都认为佟氏家族是自己人，两方都在争取他们。\NoteRef{1}{21}在%
辽东地区，佟氏是少有可算作书香门第的家族。1616年，佟卜年中了%
明朝进士。当时明朝廷经营辽东的熊廷弼就把佟卜年招致麾下，希望%
他能作为佟氏的代表，帮自己打理辽东。%
\par
% 原书页码：010；段落
佟氏家族中的佟养性和佟养真在1619年投奔到努尔哈赤帐下。努%
尔哈赤按照游牧民族的老传统，嫁了自己氏族中一位女子给佟养性，%
把他称为额驸（驸马），这样通过婚姻纽带来巩固双方的关系。\NoteRef{1}{22}在明%
廷看来，两位佟氏重要人物投靠努尔哈赤，意味着整个佟氏家族的叛%
变。因而佟氏家族的成员在明朝疆域内集体受到了牵连。按照佟家后%
人的记载，佟家成员一些被迫逃进山区，有些不得不隐姓埋名来躲避%
明廷的追查。\NoteRef{1}{23}佟卜年也被查下狱，最后冤死狱中。至死，佟卜年都%
在喊冤，死前还留下了一份绝笔痛陈其清白。\NoteRef{1}{24}%
\par
% 原书页码：010；段落
满人没有介意佟氏和明朝的种种联系，反而视佟氏家族与汉人的联%
系为佟氏的特有价值所在。努尔哈赤和后来的皇太极都信任佟氏家族的%
人，佟氏家族的人在满人军队和朝廷都得到重用。1620年代以后，皇太极%
更是直接把所有与汉人相关的事务交给佟养性来打理。这包括满人区域%
内的汉人、明疆域内迁徙过来的移民，以及在战争中抓到的汉人俘虏。\NoteRef{1}{25}%
\par
% 原书页码：010；段落
利类思、安文思在奏本中提到的佟图赖是佟养真的二儿子，是满%
人朝廷中“佟二代”的领军人物。从军事上来说，佟图赖声名显赫，%
不仅在最初满人的辽东扩张中有各种军功，在后来攻占明朝本土、消%
灭南明势力时，也立了大功（他带的部队人关以后一直打到广西）。%
人关以后不久，顺治皇帝还娶了佟图赖的女儿。1654年，也就是利%
类思、安文思上奏本的前一年，佟氏给顺治皇帝生了一个儿子。这个%
儿子就是后来的康熙皇帝。%
\par
% 原书页码：011；段落
对于佟氏家族，跟本书相关的一个问题就是，为什么佟图赖会在%
豪格死后，同意接收利类思、安文思成为自家的奴才？而且让他们二%
人吃住在自己的府上，也即让他们做的不是普通奴才，而是包衣奴才？%
“包衣”是满语的音译。据清史前辈孟森解释：“包”是满语“家”的意思，%
而“衣”相当于汉语中的虚字，类似于“之”字。包衣奴才就是指跟主%
子关系最近的、家周围的奴才。\NoteRef{1}{26}满人入关以后，包衣奴才专指满人权%
贵家里面的奴才，是整个奴才群体中地位最高的一群。后来许多清朝%
名人都是包衣奴才出身。比如曹雪芹的祖父曹寅，他就是康熙的包衣。%
\par
% 原书页码：011；段落
对于佟图赖为什么会把两位传教士收入府中做他的包衣，当时在%
北京的汤若望在写给欧洲的书信中解释过，他说是因为佟图赖和西洋%
人的关系一直都好，和他自己的私交也很好。\NoteRef{1}{27}汤若望所谓的“和西%
洋人关系好”指什么？如果真是这样，那么这个好关系又是怎么建立%
起来的？这些问题和佟氏家族及清朝汉军八旗建立的历史有关。%
% 原书页码：011；标题
\subsection{满人对西洋人的印象}
\par
% 原书页码：011–012；段落（跨页续段）
在满人最初辽东发迹的过程中，明军使用的西洋“红衣大炮”一%
直是满人骑兵的心腹大患。当明军在城池上架起西洋大炮，战无不克%
的满人骑兵基本束手无策。17世纪二三十年代，满人军队打的败仗基%
本都和明军用大炮有关。努尔哈赤自己就是在1626年的宁远战役中%
被大炮所伤，几个月后就死了。\NoteRef{1}{28}皇太极在即位后第一年，就兴兵为%
父报仇，但都败在宁远和锦州的大炮之下。不过这也让皇太极认识到%
了大炮的威力，开始致力于发展满人的炮兵部队。在1620年代后期，%
明满双方都意识到了西洋大炮在战争中起到的关键作用。1626年宁远%
大捷以后，明廷甚至派礼部官员给立功的红衣大炮做了祭祀典礼，望%
其能够继续发挥作用，以保大明江山。\NoteRef{1}{29}袁崇焕也在宁远大捷后给朝%
廷奏事时总结道，对付满人最有效的策略就是“凭坚城以用大炮”。\NoteRef{1}{30}%
同样，满人在大凌河之战中使用大炮获得胜利后，也确定行军打仗，%
大炮必不可缺。\NoteRef{1}{31}%
\par
% 原书页码：012–013；段落（跨页续段）
但在1620年代末期，明朝还没有技术仿制欧式的“红衣大炮”，%
而主要依赖从澳门的葡萄牙人那里进口，以及依靠他们训练军队，提%
高大炮的命中精度。1625年，已经受洗成为天主教徒的明朝大臣徐光%
启派了两个门人到澳门去谈购炮之事。当时在澳门的葡商以及传教士%
都认为大炮贸易是天赐良机，可以借此加强和大明的关系。所以澳门%
的商人和传教士对大炮一事都很上心。1629年，耶稣会神父陆若汉%
（João Rodrigues，1561—1634）带领 30 位葡萄牙炮手和 40 尊大炮从%
澳门出发去往北京。\NoteRef{1}{32}1630年1月，队伍在抵达距离北京约60公里%
的小城涿州时，正好遇到了满人南侵围城。他们一行便在城上架起了%
大炮迎敌。据记载，满人骑兵仅仅是听到如雷般的炮声，还没有开始%
真正攻城，就撤退了。\NoteRef{1}{33}有了这个战绩，自然整支队伍到北京后受到%
了相当高的礼遇，而且北京还给团队中每个人开出了很高的薪俸。当%
时只有带队的陆若汉因为自己的神父身份拒绝了明廷的俸禄。明军的%
前线将领得知这批带着大炮的洋人队伍来京以后，都纷纷上书朝廷要%
这些洋人和洋炮。\NoteRef{1}{34}在这一背景下，明廷在5月又派陆若汉带人到澳%
门磋商进一步购买大炮和雇佣葡萄牙炮手的事宜。11月，陆若汉带着%
100名葡萄牙炮手和12门炮从澳门启程赶赴北京。可惜的是，当时明%
廷内部廷议不能平息，一些大臣反对雇佣西洋人。整个队伍到达南昌%
以后，朝廷又命整支队伍返回澳门候旨。\NoteRef{1}{35}%
\par
% 原书页码：013；段落
在满人那边，皇太极命佟养性负责办理所有大炮事宜。1631年2%
月8日，在佟养性监督下，满人制造出了第一批大炮。其中最大的一%
尊取名为“天佑助威大将军”。佟养性的名字连带一帮负责造炮工匠%
的名字都刻在这尊大炮的炮身上。\NoteRef{1}{36}七个月以后，佟养性就用这批大%
炮在大凌河之战中立下了首功。\NoteRef{1}{37}乾隆年间，清朝官修历史总结“红%
衣大炮”的章节，写道：%
\par
% 原书页码：013；段落
{\bookkai 其年，征明，久围大凌河，而功以成用“大将军”力也，自后，}%
{\bookkai 师行必携之。}{\bookkai \NoteRef{1}{38}}%
\par
% 原书页码：013；段落
大炮生产出来的那一年就在艰难的大凌河一战中发挥了关键作用。这%
就是后来满人出征都会携带大炮的原因。%
\par
% 原书页码：013；段落
官修历史中没有记载满人的造炮技术到底是从哪里学来的。在%
1630年代，明朝这边是没有技术制造出欧式大炮的，主要依赖进口。%
当年修史的人大概也发现了这个问题，只能用春秋笔法留下疑问，没%
有给出答案：%
\par
% 原书页码：013；段落
{\bookkai 先是连鸟枪尚未造，造炮自此始。\NoteRef{1}{39}}%
\par
% 原书页码：013–014；段落（跨页续段）
鸟枪都造不出的满人，突然技术大进，造出了大炮。这个技术的跨越，%
史官只能感叹，没办法解释。清史学者牟润孙猜测，满人可能直接从%
西洋人那里得到了技术，因为负责造炮的佟养性在投奔满人之前是在%
北京做买卖的商人，那时他接触到了西洋人。\NoteRef{1}{40}尽管牟润孙没有能给%
出证据来支持这一猜测，但似乎这是唯一能够解释为什么满人突然能%
够超越明朝，自己把欧式大炮造出来的原因。\NoteRef{1}{41}%
\par
% 原书页码：014；段落
可以肯定的是最晚至1633年，满人军队中就有葡萄牙的炮手了。%
这一年，孔有德带领一支明朝部队投奔了皇太极。这支部队是当时%
明朝最精锐的炮兵部队，由受洗了天主教的孙元化一手创建。孙氏%
自己就是明朝最有名的造炮专家。\NoteRef{1}{42}在1632年冬，耶稣会神父陆若%
汉和他带进中国的葡萄牙炮手们就和这支部队一起驻扎在登州。他%
们经历了孔有德带领部队叛变的全过程。在孔有德部队围困了登州%
以后，陆若汉半夜借着大雪的掩护，跳下城墙，突围出去，通知了%
北京朝廷孔有德部队叛变的消息。\NoteRef{1}{43}葡萄牙士兵有12人在和叛军交%
战时战死。最后叛军占领登州，剩下13名葡萄牙士兵跟着叛军一起%
投降了满人。\NoteRef{1}{44}%
\par
% 原书页码：014–015；段落（跨页续段）
孔有德叛军选择走海路投降满人，这样他们可以把大炮弹药等辐%
重一起带上。当皇太极知道这支精锐部队来降之时，大喜过望，亲自%
骑马出城十里迎接。据估计，这支部队有27门当时最先进的欧式大%
炮（而当时整个满人部队只有12门），以及12000士兵外加经过葡萄%
牙炮手训练出来的明炮兵精锐。\NoteRef{1}{45}由于这支部队的特别战斗能力，皇%
太极破例没有将其编人到任何八旗中去，而是让其保留独立建制，成%
为一支劲旅。\NoteRef{1}{46}1633年8月，在投降皇太极三个月以后，这支部队就%
在攻打旅顺的战役中大放异彩。旅顺是当时辽东的一个战略要地，城%
墙坚固，装备也很完善。尽管如此，旅顺还是经不住孔有德部队大炮%
的持续攻击。历史记载，旅顺在孔有德炮兵连续七天炮击后，被满人%
士兵破城拿下。\NoteRef{1}{47}后来的战争中，孔有德部队都是满人南下的先锋，%
他们用大炮为满人骑兵拿下一个个战略要地扫清了道路。\NoteRef{1}{48}从某些方%
面来说，这支部队不仅大大提升了整个满人军队的炮兵实力，而且还%
让满人见识了西洋大炮和西洋人训练出来的炮兵技术。%
\par
% 原书页码：015；段落
1631年大凌河一役后，皇太极不仅认识到了大炮的作用，同时也%
认识到了汉人在军队中的特殊作用，因为制造和使用大炮这一新式武%
器都离不开汉人。在满人社会中，以及在明朝北方的所有游牧民族中，%
军功都是衡量社会地位的主要标准。当汉人使用的西洋大炮在战争中%
发挥作用以后，汉人在军事上的实力也大大提升了，发挥了满人、蒙%
古人都无法替代的作用。在这个背景下，皇太极在大凌河胜利后，开%
始把汉人归结到一起，组建独立于满旗和蒙古旗以外的汉军旗编制。\NoteRef{1}{49}%
成立汉军旗也便是正式确认了大炮部队成为满人军队的特别兵种。这%
样1620年代以来形成的明朝炮兵与满人骑兵之间的军事平衡就被打%
破了。满人在其骑兵基础上建制了炮兵，军事的天平就开始向满人军%
队一侧倾斜了。对于这一点，史家陈寅恪在读这段历史的时候，称这%
是“明清兴亡之一大关键”。\NoteRef{1}{50}%
\par
% 原书页码：015–016；段落（跨页续段）
至少，满人越来越多地使用大炮，使清朝疆域内汉人士兵成为满%
人部队的重要部分，同时提升了汉人整体的社会地位。乾隆朝官修%
的《皇清开国方略》中就坦陈红衣大炮的使用是成立汉军旗的主因。%
同样也是这个原因，负责制造大炮的佟养性成为了第一支汉军旗的固%
山。当时皇太极在任命的时候，指出佟养性负责管理所有跟汉人有关%
的事务。同时皇太极还预感到，从各旗中抽调出汉人来集中管理，成%
立新的旗，肯定会损害到一些王公们的利益，所以还专门强调在汉%
人相关的问题上必须要服从佟养性的安排。\NoteRef{1}{51}对于集中在一起单独成%
旗的汉人来说，新成立的汉军旗无疑大大提高了汉人在满人社会中%
的地位。\NoteRef{1}{52}%
\par
% 原书页码：016；段落
第一支汉军旗成立于1631年，1637年发展成两旗，1639年变成%
四旗，最后1642年发展成了八旗。\NoteRef{1}{53}每一次汉军旗的扩展都和满人缴%
获大量火器和汉人俘虏有关。1637年3月，满人终于在占领朝鲜后，%
拿下了皮岛。皮岛，现称椴岛，属于朝鲜。地理上看，该岛离朝鲜更近。%
1620年代以来，毛文龙率军驻扎在皮岛，一直是满人南下时后方的一%
块心病。该岛士兵以使用欧式红衣大炮闻名于世。该岛的炮兵是葡萄%
牙人训练出来的。1631年皇太极御驾亲征皮岛之时，守岛的炮兵就是%
在葡萄牙人带领下击退了满人的进攻。1637年的登岛之战是鳌拜的成%
名之战，他带领一支敢死队突破了守军防线，成功拿下了三个月都没%
有攻下的皮岛。登岛后，按照惯例，作为不投降的惩罚，满人屠杀了%
全部守军。所以皮岛之战的胜利，基本没有增加多少汉人俘虏，满人%
主要得到的是皮岛上面大量的火器和弹药（满人军队在岛上缴获了约%
10门红衣大炮，而当时满人已有约38门）。\NoteRef{1}{54}在1637年9月，也就%
是攻占皮岛后半年，汉军旗从一旗发展成了两旗。%
\par
% 原书页码：016–017；段落（跨页续段）
1638年秋，皇太极开始从两个方向进攻明朝。1639年，满人军%
队占领山东大部分地区，在济南缴获大量火器和弹药。山东人口密集，%
当时占领区有超过50万的人口。为了管理大量新缴获的火器以及汉%
人，两旗的汉军增加到四旗。1642年，在取得松锦之战胜利后，满人%
又获得了大量的火器和汉人，汉军旗进一步扩展到了八个旗。这样汉%
军八旗就和满八旗以及蒙古八旗形成了相等的建制规模。至此，满人%
军队据称拥有了一百门红衣大炮，而明朝只剩下10门红衣大炮可用%
来拱卫京师。对于满人炮兵实力的迅速增加，辽东巡抚黎玉田在1643%
年3月上奏北京朝廷，分析双方的军事力量：%
\par
% 原书页码：017；段落
{\bookkai 我之所以制酋者，向惟火器为先，盖因我有而酋无，故足以取胜。}%
{\bookkai 后来，酋虽有而我独多，犹足以侥幸也。今据回乡称说，酋于锦州铸}%
{\bookkai 造西洋大炮一百位，我之所有曾不及十分之一。设不幸，卒如回乡所}%
{\bookkai 言，酋以大炮百位排设而击，即铁壁铜墙亦恐难保。\NoteRef{1}{55}}%
\par
% 原书页码：017；段落
这一前线的报告对于双方军力消长的分析，就是立足于火器对比之%
上的。应该说双方都深知大炮对于整个战况的影响，这份报告也从%
侧面说明了为什么善于骑射的满人从1630年代以后一直重视炮兵的%
发展。%
\par
% 原书页码：017；段落
八旗汉军旗的成立和大炮使用之间的联系还体现在汉军旗的满语%
名字上。在满语中，汉军旗叫 Ujen cooha，在汉文中直接音译为“乌%
珍超哈”。\NoteRef{1}{56}在满语中，汉人叫nikan，写成汉字音译是“尼堪”。这%
里就有一个问题，为什么汉军旗不叫“nikan cooha”（尼堪超哈），而%
是用“Ujen”（本意是“重的”）这个满语词？\NoteRef{1}{57}陈寅恪在《柳如是别传》%
中已经注意到这个问题，并有简要论述：%
\par
% 原书页码：017–018；段落（跨页续段）
{\bookkai 至满洲语所以称“汉军”为“乌珍超哈”，而不称为“尼堪超哈”}%
{\bookkai 者，推原其故，盖清初夺取明室守御辽东边城之仿制西洋火炮，并用}%
{\bookkai 降将管领使用，所以有此名号。\NoteRef{1}{58}}%
\par
% 原书页码：018；段落
关于“乌珍超哈”和大炮的联系，皇太极在成立第一支汉军旗时做的%
安排就揭示了其中的奥秘，%
\par
% 原书页码：018；段落
{\bookkai 集诸贝勒大臣……其随营红衣炮、大将军炮四十位，及应用挽车}%
{\bookkai 牛骡，皆令总兵官佟养性管理。\NoteRef{1}{59}}%
\par
% 原书页码：018；段落
用牛车和骡车拉着的大炮，对于满人骑兵来说，肯定是“重”装备。%
当皇太极下令把这些重型器械收归到一处，统一交给汉军旗管理以后，%
各支满蒙旗都如释重负，这样他们便可以更好地发挥他们快速奔袭的%
骑兵特色。%
\par
% 原书页码：018；段落
就像是在经济发展中出现的社会分工一样，汉军旗的成立其实是%
满人根据战争需要而进行的一次部队的军事分工。1633年，皇太极把%
满人权贵聚集到一起，研究攻打朝鲜和大明的方略。豪格就提出应该%
让满蒙骑兵从北京的北面南下，穿过大漠，闪击北京。同时让汉军从%
北京的东北方向携带大炮一个城池一个城池地进攻。\NoteRef{1}{60}豪格的这个建%
议，就很明显地表明满人军队在创立汉军旗的时候，对汉军旗携带大%
炮的战术特点有清晰的认识。%
\par
% 原书页码：018–019；段落（跨页续段）
另外，满人使用“重”（Ujen）而不是用“汉”这个民族概念来命%
名“汉军”，还有可能是为了避免不必要的民族矛盾。因为当时在八旗%
的“汉军”中有很多使用汉语，但不是民族意义上的“汉人”，最有名%
的就是契丹人的后代，有些书中称为辽人，他们在东北的满人领地内%
属于汉化程度很高的一个族群。另外佟养性之后的第二位汉军旗统领%
石廷柱，就肯定不是汉人（有可能是契丹人的后代，也有说是满人）。%
据《清史稿》记载，他家从他父亲一辈才开始使用汉族的“石”姓。\NoteRef{1}{61}%
% 原书页码：019；标题
\subsection{神父与大炮生意}
\par
% 原书页码：019；段落
耶稣会神父自明末入华以后，他们在中国的发展与明清双方引进%
西洋大炮都有关联。他们希望借助大炮贸易进一步接触到明清双方的%
权贵阶层，以便推进他们的传教事业。在1630年成功和明廷做了第%
一笔大炮生意以后，耶稣会神父陆若汉在第二批运送大炮的队伍中加%
了五位神父进去。由于第一笔生意受到了明廷的肯定，因而第二支队%
伍从澳门出发进人内地后，一直受到沿途官员的优待。\NoteRef{1}{62}尽管后来明%
廷内党派之间不能统一意见，不能决定到底是不是应该信任洋人，在%
这支队伍走到南昌以后，让他们原路返回了，但陆若汉安插的五位神%
父还是悄悄地按照原定计划去到了各自的目的地。\NoteRef{1}{63}同时，大名鼎鼎%
的汤若望神父，其实除了在明廷做一个天文历法的官员外，还在教授%
造炮技术。\NoteRef{1}{64}他的讲解在明末编辑成书，名为《火攻挈要》。\NoteRef{1}{65}%
\par
% 原书页码：019–020；段落（跨页续段）
有了大炮这个引子，才能了解传教士和八旗中汉军旗之间的关系。%
在满人中，使用西洋大炮的汉军就是八旗中最容易接纳西洋传教士的%
人群。最明显的一个例子就是孔有德解救耶稣会神父谢务禄（Alvaro%
\space%
Semedo，1585—1658，又名曾德昭，葡萄牙人）。与汤若望神父投降%
清廷不一样，谢务禄一直死忠于明廷，后来追随南明一直南撤，直到%
在广州被清朝士兵俘虏。在清朝士兵要对谢务禄行刑之前，孔有德得%
知他是西洋传教士，亲自出面救下他。孔有德还把他接到自己的住处，%
礼遇有加，并说降了谢务禄，而后委派他作为使节到澳门去和葡萄牙%
人议和，因为葡萄牙人不久前刚派了一支军队到广州。\NoteRef{1}{66}%
\par
% 原书页码：020；段落
孔有德是当时最为著名的降清汉人之一，他带着部队一路从东北%
打到南方。心怀大明的汉人都恨他入骨，认为他是最大的汉奸。问题%
是像他这样的清朝汉军统帅，为什么会出手救一位素昧平生的耶稣会%
神父？这就是因为耶稣会神父与当时的大炮军火购买有密切关系。当%
孔有德还是明朝将领的时候，他的顶头上司是孙元化。孙元化是明末%
受洗天主教的著名高官之一，同时也是明朝中大炮火器的专家。孙元%
化的部队长期以来都和葡萄牙人以及耶稣会神父往来密切。当孔有德%
率部队叛变准备降清的时候，部队中就有三十几位葡萄牙大炮教官以%
及一名耶稣会神父。\NoteRef{1}{67}所以降清以后，孔有德部队一直是清朝军中使%
用火炮的王牌军。孔有德和耶稣会神父的渊源可以说就是来自于他的%
部队对大炮军火的使用。通过军火交易，孔有德接触过耶稣会神父，%
他深知这些神父的用处。这就是为什么他会出手搭救一位自己不认识%
的神父，而且还能很快说服他，让他在与葡萄牙的谈判中为自己效力。%
对孔有德因为火炮的因缘搭救谢务禄，当时传教士也都知道，也告知%
了欧洲。\NoteRef{1}{68}%
\par
% 原书页码：020–021；段落（跨页续段）
另外，明亡以前到中国的传教士基本都只学过汉语，不懂满语。%
这样八旗中的汉军自然而然就成为传教士的依靠。1646年，利类思和%
安文思神父在四川遭遇豪格大军，下到士兵上到豪格，都不会汉语。%
多亏八旗汉军的一位随军翻译从中斡旋，才使二位神父得到豪格的优%
待。巧合的是，在1644年满人人京的时候，汤若望神父也和这位翻%
译打过交道。\NoteRef{1}{69}这位翻译的名字背景已不可考，但同一个人两次出现%
在传教士面前，已经说明清朝八旗中的汉军是传教士在明清易代之际%
的贵人。\NoteRef{1}{70}%
% 原书页码：021；标题
\subsection{佟氏家族与传教士}
\par
% 原书页码：021；段落
耶稣会神父具体是什么时候和佟氏成员开始往来联系的，现在%
已不可考。但有详细记载的是，在满人军队人京的第二天，汤若望%
神父上了一道陈情书给范文程。范文程是当时清朝负责恢复北京战%
后秩序的核心官员。满人进京后，第一项大工程就是要把北京城北%
面的住户全部迁到城南，城北归满人贵族特有。这样的大工程就涉%
及大量住在城北的达官贵人，这些人都希望得到范文程的特许，允%
许他们不迁出城北。汤若望的教堂就在城北，他也用中文写了陈情%
书，希望他的教堂得到特许。根据汤若望自己的记载，由于他是外%
国人，他写的条陈才得以送到了范文程手中。范文程后来派了几个%
人去汤若望的教堂实地考察了一番，最后特许教堂不用搬迁。\NoteRef{1}{71}%
\par
% 原书页码：021；段落
从家世来说，范文程和佟氏渊源深厚。他的姐姐就嫁到了佟家。%
由于佟氏是名门望族，部分佟氏名人降清以后，明廷震怒，在明朝土%
地上的佟氏族人受到迫害，范文程曾通过各种渠道出手相助。佟世思%
（1652—1692）在其写的家传《先高曾祖三世行略》中，就记载了佟%
氏受到迫害后，范文程对佟氏的种种保护。满人人关以后，范又帮助%
佟氏谋取官职。\NoteRef{1}{72}%
\par
% 原书页码：021–022；段落（跨页续段）
范文程自称是范仲淹的17代孙。1619年，他投奔努尔哈赤，显%
赫的家世背景使得努尔哈赤重用了他。\NoteRef{1}{73}也是这一年，佟氏家族的佟%
养性和佟养真投奔了努尔哈赤。皇太极1636年在沈阳称帝后，要按%
照明朝廷的规模来积极组建自己的朝廷和汉式的官僚体制，范文程是%
负责人。当范文程得知李自成攻陷北京后，认为是天赐良机，上书催%
促多尔衮出兵，这大概是他对满人最大的贡献。\NoteRef{1}{74}人京以后，他负责%
维系保证已有的明朝官僚体系继续运作。%
\par
% 原书页码：022；段落
范文程熟知中国官僚系统运作方式，因而他知道精通天文历法的%
汤若望神父对清初稳定天下有大用处。天文历法跟中国传统中的“天%
命”观念紧密相连。比如，1644年9月1号，出现了日食，与汤若望%
测算的误差仅有几分钟；汉官的测算误差是半小时，而回回的测算误%
差为一小时。\NoteRef{1}{75}范文程明白，如果大清公布的历法能切合天象，就能%
让老百姓在心理上接受清取代明是天命使然。因此，范文程不仅特许%
汤若望的教堂不用搬迁，还让他掌管钦天监，成为朝廷命官。汤若望%
后来和范家保持了密切来往，范文程不仅给汤的教堂捐了钱，他的女%
儿以及一些家丁还都信了教。多尔衮去世以后，范文程还动用私人关%
系把汤若望引荐给了顺治皇帝。\NoteRef{1}{76}%
\par
% 原书页码：022–023；段落（跨页续段）
在满人人京的第一个月中，代善（1583—1648）多次去教堂拜访%
汤若望，谈论满人历史。\NoteRef{1}{77}代善是努尔哈赤的二儿子，自1615年努尔%
哈赤的大儿子去世后，代善就是努尔哈赤家族中的长子。在努尔哈赤%
和皇太极去世后的两次汗（帝）位争夺中，代善通过协调各方，避免%
了满人内部的分裂。在1644年的北京，代善理所当然地是当时最重%
要的亲王。这也是为什么汤若望特意把代善来访的消息记录下来，写%
在邮寄回欧洲的信中。现在已不可考到底是谁把汤若望引荐给代善的，%
但可以确定的一点是：代善母亲所属的满人佟佳氏和康熙母亲的佟氏%
家族有很深的渊源。佟家后来也改姓为佟佳。%
\par
% 原书页码：023；段落
京城以外，明清交替之际，佟国器在江南地区关照传教士是当时%
天下皆知的。佟国器的父亲佟卜年是明末进士，还在明廷做过官。后%
来受佟氏家族降清的影响，佟卜年下狱并冤死狱中。佟卜年死后，佟%
国器和他的母亲得到当时佟卜年江南师友的资助，辗转住在武昌、南%
京，最后定居宁波。佟国器的母亲陈氏出身书香门第，在佟卜年中进%
士后在抚顺等候官缺之时嫁入佟家。据陈寅恪考证，陈氏诗文水平颇%
高，能和江南才子们唱和往来。\NoteRef{1}{78}佟国器除了父母都是能文之人，自%
身也是典型的传统文人，也有诗文传世。从文化角度看，佟国器是典%
型的江南士人。但佟国器和江南士人最大的不同就是他的籍贯，他属%
于抚顺的佟氏家族。正是因此，当佟图赖1645年率领的正蓝旗军队%
打到嘉兴以后，便把佟国器招至旗下，使他从一介平民一跃成为地方%
官，管理江南最富裕的浙江嘉兴湖州片区（浙江嘉湖道）。\NoteRef{1}{79}%
\par
% 原书页码：023–024；段落（跨页续段）
作为佟氏家族在江南的代表，佟国器仕途亨通，一路青云直上。%
1653年升任福建巡抚，1658年调任浙江巡抚。佟国器一家和传教士%
有很深的渊源。他的夫人就是受洗的天主教徒，在西文中的天主教名%
字是Agathe。1645年他出仕为官以后，就一直在帮助传教士。在明清%
之交的动荡年代，意大利卫匡国（MartinoMartini）神父在江南的活动%
就多仰仗佟国器的关照。在江南为官期间，佟国器还在多地为天主教%
重修教堂。\NoteRef{1}{80}他跟当时很多神父都有往来，他还给神父阳玛诺（Manuel%
\space%
Diaz,1574—1659）、何大化（Antonio de Gouvea,1592—1677）和贾谊%
睦（Girolamo de Gravina, 1603—1662）出版的神学著作写了序言。\NoteRef{1}{81}简%
单来说，佟国器支持保护天主教在当时是很出名的。后来，在1664年%
反天主教的教案中，他被归为天主教“三大护法”之一。\NoteRef{1}{82}%
% 原书页码：024；标题
\subsection{消除奴籍}
\par
% 原书页码：024；段落
传教士通过西洋大炮火器和八旗中的汉军搭上了线，和汉军旗中%
负责大炮的佟氏家族有渊源往来。明白了这个基本背景之后，才能理%
解为什么豪格去世以后，利类思和安文思两位神父的奴籍会过户到佟%
图赖家去。而且，当汤若望知道两位神父过户到佟图赖名下以后，大%
赞这是最好的结果。%
\par
% 原书页码：024；段落
从佟图赖优待两位神父的种种具体表现来看，他无疑是亲天主教%
的八旗权贵之一。佟图赖把两位神父的吃住都安排在了自己府上。\NoteRef{1}{83}%
也就是这段时间，两位神父和佟图赖的两个儿子佟国维、佟国纲建立%
了深厚关系。他们二人和佟图赖家族的关系不是什么秘密，无论在当%
时的中国还是欧洲，都是广为人知的。就是七十年以后的1723年，巴%
多明神父在一封信中写到和雍正朝隆科多的交往时，还在念叨佟家和%
传教士的长期关系，还提到了这两位神父住在佟家的往事。\NoteRef{1}{84} 隆科多%
就是佟国维的儿子，他和传教士的故事在后面会讲到。%
\par
% 原书页码：024–025；段落（跨页续段）
1653年，佟图赖帮两位神父消除了奴籍，让二人成为自由身。佟%
图赖具体怎样操作解除奴籍的细节已不可考。但是据传教士留下的记%
录，可以看到佟图赖为两位神父消除奴籍没少费周章。他曾经几次摆%
席宴请官员来打通关节。\NoteRef{1}{85}两年之后，也就是1655年，在佟家的引荐%
下，顺治皇帝赏钱赐地给两位神父，帮助他们修了属于自己的教堂。\NoteRef{1}{86}%
对于两位神父来说，佟图赖作为主子能帮他们消除奴籍是对他们天大%
的恩赐。\NoteRef{1}{87}在满人的传统中，只有对主子家庭有大贡献或者在战场上%
立了大功，才有可能在主子和朝廷相关部门都同意的情况下，解除奴%
籍。\NoteRef{1}{88}对于两位神父来说，去除了奴籍才能更好地、不受约束地传教，%
接触信众。但同时，根据满人传统，解除奴籍以后，奴才和前主子之%
间的派系、政治纽带关系并不会解除。\NoteRef{1}{89}这样一来，在朝堂上，两位%
神父还是佟家一系。明白了他们和佟家的这层关系，也才能明白为什%
么在本篇开头的奏本中，他们要特别提到佟图赖。因为两位神父明白，%
在清初的朝堂政治中，家族派系是一切活动的基础，而佟家就是他们%
在满人世界中的“娘家”。%
% 原书页码：025；标题
\subsection{小结}
\par
% 原书页码：025；段落
1655年，也就是利类思和安文思两位神父递上开篇那份奏本的当%
年，他们盼望已久的教堂终于落成。教堂取名圣若瑟教堂。由于其位%
置在紫禁城东边，又俗称东堂。这座教堂今天是北京的一个景点，在%
现在的王府井大街74号。这座教堂凝聚了两位神父多年的心血。%
\par
% 原书页码：025–026；段落（跨页续段）
长期以来，传教士的历史都忽略了利类思和安文思二位神父的存%
在。历史书写都把目光集中在了传奇人物汤若望身上。汤公在满人人%
关第二天就获礼遇，被邀入朝，成为一个典型的中式大臣。再加上汤%
公曾在明末朝廷做官，李自成入京以后也对他礼遇有加，这些经历注%
定他是基督教历史中的传奇人物。汤公在清廷做官时，言谈和书写都%
使用汉文，而且顺治皇帝还经常邀他入宫深谈。作为第一个在顺治皇%
帝身边的传教士，汤若望很自然地被视为传教士步人清宫的源头，他%
和顺治的君臣关系也被视为后来康熙与传教士关系的源头。%
\par
% 原书页码：026–027；段落（跨页续段）
由于利类思和安文思1648年到北京后，汤若望拒绝他们住在%
他的教堂中，二位神父就一直希望有一座属于他们自己的教堂。据%
汤若望的记载，二人被俘人京后，他的朝廷眼线告知他二人的身份%
是张献忠朝廷的罪臣，而且已归为奴籍，让他不要贸然接触。同%
时，汤若望也请示过他所在教区的上级傅汎际（FranciscoFurtado，%
1589—1653），傅也指示不要接触，以免影响到整个北京的传教事%
业。\NoteRef{1}{90}龙华民（Nicholas Longobardi，1559—1654）也是同样的意见：%
不要接触，以免影响整个中国的传教事业。利玛窦以后，龙华民接%
任中国耶稣会的会长。他资历老，此时已经九十多岁。他在1651年%
写过一封信给罗马，解释利类思和安文思二人为奴的情况，也指出%
了收留二人有可能给整个传教事业带来麻烦。\NoteRef{1}{91}后来据写《汤若望传》%
的魏特考证汤若望的各种书信，认为他其实一直关注着二人的案子，%
在可能的情况下暗暗帮助他们。不过，汤若望和这两位神父还是结%
下了怨。1650年之后的十来年中，两位神父一直在向罗马教廷控告%
汤若望，指控他作为神父人朝为官、参与算命，以及私生活不检点。%
中国大部分的传教士也都被牵涉进来，接受罗马来来回回的各种调%
查问讯。\NoteRef{1}{92}在经过数十年的调查后，最后罗马裁决汤若望无罪。\NoteRef{1}{93}一%
直到1664年，也就是汤若望去世前一年半，两位神父才和汤若望改%
善了关系。因为这时清初著名的历狱开始了，而案子针对的是整个%
天主教群体，他们不得不和汤若望站在一起，共同应对。从1648到%
1664年，两位神父和汤若望神父之间有16年时间的对立，这也意味%
着两个群体在满人占领下的北京按着自己的路线互不来往地开辟着%
各自的传教事业。%
\par
% 原书页码：027；段落
汤若望在满人人京后降清，然后仕途一帆风顺，成为了清初大臣，%
执掌钦天监，这些都是历史中为人熟知的。然而利类思和安文思二位%
神父是在战场上与满军相遇，在战斗中，利类思被箭射中大腿，安文%
思被射中手臂。\NoteRef{1}{94}在战场上的生死瞬间，满人士兵因为发现二人的西%
洋人长相，竟然在手起刀落的一瞬，突然停下了砍向两位神父的刀。%
两位神父用文字记录下了这生死瞬间，并为此感谢上帝，认为是上帝%
救了他们。\NoteRef{1}{95}不过，对于读史的人来说，一定有疑问，为什么满人士%
兵会有这样的意识，认为西洋人有用，从而刀下留人？%
\par
% 原书页码：027；段落
对于明清之交的满人士兵来说，他们所能直接熟知的西洋利器%
就是大炮。早在1633年，豪格论述攻明良策的时候，就坦言用军中%
的汉军炮击攻城才是正道，而他所掌管的正蓝旗下的汉军旗一直以%
火炮攻击见长。\NoteRef{1}{96}所以当豪格军中士兵在1647年的战场上遇见二位%
神父，看见他们的西洋面孔，他们就知道这二人可能有用。西洋大%
炮的威力，大炮对战事的作用，跟着豪格一路从东北打到四川的满%
人士兵是知道的。因而与其说是杀红眼的满人士兵突然感悟到了上%
帝之光而放下了屠刀，不如说是西洋大炮的威力救了两位神父的命。%
大炮和西洋面孔的绑定思考并不是猜想，而是当时亚洲世界的普遍%
情况。有在印度的传教士就提到无论是否参与过大炮贸易，有些地%
方规定，所有西洋面孔的人都不能通行，原因就是西洋大炮在那些%
区域的使用。\NoteRef{1}{97}%
\par
% 原书页码：027–028；段落（跨页续段）
二位神父在战场上捡回性命后，被归入奴籍，挂入了清朝八旗中。%
他们归顺清朝统治的轨迹和汤若望神父完全不一样。汤若望跟许多降%
清的汉臣轨迹相似，改朝换代后又直接成了新朝廷的大臣。而二位神%
父是以战俘奴隶身份进人到满人世界中的。他们的这一历史轨迹在后%
来的历史中已经被慢慢淡忘了。然而他们二人才是对天主教后来在华%
发展真正产生了决定性影响的人。特别是他们的奴籍在豪格死后转到%
了佟图赖名下，而佟图赖又是康熙皇帝的外公。也就是说，由于佟图%
赖的关系，两位神父被归人了康熙的外戚圈子之中。这一层关系，为%
后来传教士进入康熙的世界开启了大门。这其中的故事就是下面三章%
的内容。%
